\chapter{Toy Model: One-Loop Effective Potential and Renormalization}
\label{ch:toy-model}

We illustrate the general formalism of the previous chapter with an
explicit computation in the simplest setting: a single real scalar field
with Lagrangian
\begin{equation}
    \mathcal{L}(x) = \frac{1}{2}\,\partial_\mu\phi(x)\,\partial^\mu\phi(x)
    - V(\phi(x)),
    \label{eq:real_scalar_field}
\end{equation}
and classical potential
\begin{equation}
    V(\phi) = \frac{\alpha_0}{4!}\,\phi^4 - \frac{m_0^2}{2}\,\phi^2.
\end{equation}
We work in units $\hbar = c = 1$ throughout.

\section{The Classical Field and the Effective Action}

The classical minimum of $V(\phi)$ sits at some $\phi_0$, but the vacuum
expectation value $\langle\phi\rangle$ of the quantum field need not
coincide with it. It is defined by
\begin{equation}
    \langle\phi\rangle \equiv \lim_{J\to 0}
    \frac{\langle 0^+|\phi_{\text{op}}|0^-\rangle_J}
         {\langle 0^+|0^-\rangle_J},
\end{equation}
where the generating functional $W[J]$ encodes the full quantum theory via
\begin{equation}
    \exp\frac{i}{\hbar}W[J]
    = \mathcal{N}\int(\mathcal{D}\phi)
      \exp\frac{i}{\hbar}\left\{S[\phi] + (J,\phi)\right\},
\end{equation}
with $S[\phi] = \int d^4x\,\mathcal{L}(x)$ and
$(J,\phi) = \int d^4x\,J(x)\,\phi(x)$.
The classical field, the VEV in the presence of an external source $J(x)$,
is
\begin{equation}
    \phi_c(x) \equiv
    \frac{\langle 0^+|\phi_{\text{op}}(x)|0^-\rangle_J}
         {\langle 0^+|0^-\rangle_J}
    = \frac{\delta W[J]}{\delta J(x)},
    \label{eq:phi_c}
\end{equation}
and $\langle\phi\rangle$ is recovered as $\lim_{J\to 0}\phi_c(x)$. The
loop expansion is an expansion in powers of $\hbar\,c\,\alpha$, but it
performs an infinite resummation of Feynman graphs rather than a
perturbative truncation.

The natural question is: given a desired profile $\phi_c$, what source
$J$ produces it? To answer this, we trade $J$ for $\phi_c$ as the
independent variable via the Legendre transform,
\begin{equation}
    \Gamma[\phi_c] \equiv W[J] - (J,\phi_c),
    \label{eq:Gamma_def}
\end{equation}
where $J$ is eliminated in favour of $\phi_c$ by inverting
\eqref{eq:phi_c}. The functional $\Gamma[\phi_c]$ is the effective action,
so named because it is a functional of the classical field $\phi_c$ and
thus plays the same role for the quantum theory as $S[\phi]$ does
classically. Differentiating \eqref{eq:Gamma_def} and using
\eqref{eq:phi_c},
\begin{align}
    \frac{\delta\Gamma[\phi_c]}{\delta\phi_c(x)}
    &= \frac{\delta W[J]}{\delta\phi_c(x)} - J(x)
    - \int d^4y\,\frac{\delta J(y)}{\delta\phi_c(x)}\,\phi_c(y), \notag\\
    \frac{\delta W[J]}{\delta\phi_c(x)}
    &= \int d^4y\,\frac{\delta W[J]}{\delta J(y)}
       \frac{\delta J(y)}{\delta\phi_c(x)}
     = \int d^4y\,\frac{\delta J(y)}{\delta\phi_c(x)}\,\phi_c(y),
\end{align}
so that
\begin{equation}
    \frac{\delta\Gamma[\phi_c]}{\delta\phi_c(x)} = -J(x).
    \label{eq:eff_action_eom}
\end{equation}
At $J=0$, translational invariance forces $\phi_c$ to be a constant, so
$\langle\phi\rangle$ is the root of
\begin{equation}
    \frac{d\Gamma[\phi_c]}{d\phi_c}\bigg|_{\phi_c=\langle\phi\rangle} = 0.
\end{equation}

\section{The Effective Action as a Generator of 1PI Green's Functions}

Expanding $\Gamma[\phi_c]$ in powers of $\phi_c$,
\begin{equation}
    \Gamma[\phi_c] = \sum_{n=0}^{\infty} \frac{1}{n!}
    \int d^4x_1\cdots d^4x_n\,
    \Gamma_n(x_1,\ldots,x_n)\,\phi_c(x_1)\cdots\phi_c(x_n),
    \label{eq:gamma_expansion}
\end{equation}
we claim that $\Gamma_n$ is the $n$-point 1PI Green function, the sum
of all connected 1PI Feynman graphs with $n$ external legs. To see this,
consider the functional integral
\begin{equation}
    \mathcal{N}\int(\mathcal{D}\phi)
    \exp\frac{i}{\hbar}\left\{\Gamma[\phi] + (J,\phi)\right\}.
\end{equation}
In the limit $\hbar\to 0$ the integrand is dominated by its saddle point,
which occurs at $\phi = \phi_c$ since \eqref{eq:eff_action_eom} is
precisely the saddle-point condition. Thus
\begin{equation}
    \mathcal{N}\int(\mathcal{D}\phi)
    \exp\frac{i}{\hbar}\left\{\Gamma[\phi] + (J,\phi)\right\}
    \xrightarrow{\hbar\to 0}
    \mathcal{N}\exp\frac{i}{\hbar}
    \left\{\Gamma[\phi_c] + (J,\phi_c)\right\},
\end{equation}
and by \eqref{eq:Gamma_def} the right-hand side equals
$\exp\frac{i}{\hbar}W[J]$, so
\begin{equation}
    \exp\frac{i}{\hbar}W[J]
    \xrightarrow{\hbar\to 0}
    \mathcal{N}\int(\mathcal{D}\phi)
    \exp\frac{i}{\hbar}\left\{\Gamma[\phi] + (J,\phi)\right\}.
\end{equation}
In this limit the right-hand side is the sum of all tree graphs of a
theory whose action is $\Gamma[\phi]$, with vertices given by the
$\Gamma_n$ of \eqref{eq:gamma_expansion}. Since $W[J]$ generates all
connected Green functions of the original theory, and any connected
graph decomposes into 1PI components, it follows that $\Gamma_n$ is a
1PI Green function of the original theory. \qed

\section{The Effective Potential}

Alternatively, $\Gamma[\phi_c]$ may be expanded in derivatives of
$\phi_c$,
\begin{equation}
    \Gamma[\phi_c] = \int d^4x\left[
        -U(\phi_c(x))
        + \frac{1}{2}\,\partial_\mu\phi_c(x)\,\partial^\mu\phi_c(x)\,
          Z(\phi_c(x))
        + \cdots\right],
\end{equation}
where $U(\phi_c)$ is the effective potential. For a constant configuration
$\phi_c(x) = \rho$ this reduces to
\begin{equation}
    \Gamma[\rho] = -\Omega\,U(\rho),
\end{equation}
where $\Omega$ is the total spacetime volume. To connect $U(\rho)$ to the
1PI Green functions, we pass to momentum space. The Fourier transforms
of $\phi_c$ and $\Gamma_n$ are defined by
\begin{align}
    \tilde{\phi}_c(k) &= \int d^4x\,e^{ik\cdot x}\,\phi_c(x), \\
    \phi_c(x) &= \int\frac{d^4k}{(2\pi)^4}\,e^{-ik\cdot x}\,\tilde{\phi}_c(k),
\end{align}
\begin{equation}
    \tilde{\Gamma}_n(k_1,\ldots,k_n)\,(2\pi)^4\,\delta^4(k_1+\cdots+k_n)
    = \int d^4x_1\cdots d^4x_n\,\Gamma_n(x_1,\ldots,x_n)\,
      e^{i(k_1 x_1+\cdots+k_n x_n)}.
\end{equation}
Substituting these into \eqref{eq:gamma_expansion},
\begin{multline}
    \Gamma[\phi_c]
    = \sum_{n=0}^{\infty}\frac{1}{n!}
      \int d^4x_1\cdots d^4x_n\,
      \Gamma_n(x_1,\ldots,x_n)\,\phi_c(x_1)\cdots\phi_c(x_n) \\
    = \sum_{n=0}^{\infty}\frac{1}{n!}
      \int\frac{d^4k_1}{(2\pi)^4}\cdots\frac{d^4k_n}{(2\pi)^4}
      \int d^4x_1\cdots d^4x_n\,
      \Gamma_n(x_1,\ldots,x_n)\,
      e^{-i(k_1 x_1+\cdots+k_n x_n)}
      \,\tilde{\phi}(k_1)\cdots\tilde{\phi}(k_n) \\
    = \sum_{n=0}^{\infty}\frac{1}{n!}
      \int\frac{d^4k_1}{(2\pi)^4}\cdots\frac{d^4k_n}{(2\pi)^4}\,
      (2\pi)^4\,\delta^4(k_1+\cdots+k_n)\,
      \tilde{\Gamma}_n(k_1,\ldots,k_n)\,
      \tilde{\phi}(k_1)\cdots\tilde{\phi}(k_n).
      \label{eq:rewriting_Gamma}
\end{multline}
Setting $\phi_c(x) = \rho$ gives $\tilde{\phi}_c(k) =
(2\pi)^4\delta^4(k)\,\rho$, so each momentum integral collapses,
\begin{multline}
    \Gamma[\rho]
    = \sum_{n=0}^{\infty}\frac{1}{n!}
      \int\frac{d^4k_1}{(2\pi)^4}\cdots\frac{d^4k_n}{(2\pi)^4}\,
      (2\pi)^4\,\delta^4(k_1+\cdots+k_n)\,
      \tilde{\Gamma}_n(k_1,\ldots,k_n) \\
      \cdot\,(2\pi)^4\delta^4(k_1)\,\rho\cdots(2\pi)^4\delta^4(k_n)\,\rho \\
    = \sum_{n=0}^{\infty}\frac{1}{n!}
      \underbrace{(2\pi)^4\delta^4(0)}_{\Omega}\,
      \rho^n\,\tilde{\Gamma}_n(0,\ldots,0),
\end{multline}
and therefore
\begin{equation}
    \Gamma[\rho] = \Omega\sum_{n=0}^{\infty}
    \frac{\rho^n}{n!}\,\tilde{\Gamma}_n(0),
    \label{eq:Gamma_rho}
\end{equation}
where $\tilde{\Gamma}_n(0)\equiv\tilde{\Gamma}_n(0,\ldots,0)$ and
$\Omega = (2\pi)^4\delta^4(0)$ is the total spacetime volume. Comparing
\eqref{eq:Gamma_rho} with $\Gamma[\rho]=-\Omega\,U(\rho)$,
\begin{equation}
    U(\rho) = -\sum_{n=0}^{\infty}\frac{\rho^n}{n!}\,\tilde{\Gamma}_n(0),
\end{equation}
which is simply the statement that the Taylor coefficients of $U$ are the
zero-momentum 1PI Green functions,
\begin{equation}
    -\frac{d^n U(\rho)}{d\rho^n}\bigg|_{\rho=0} = \tilde{\Gamma}_n(0).
\end{equation}
The effective potential is thus entirely determined by the 1PI Green
functions of the original theory evaluated at vanishing external momenta.
In the next section we compute it explicitly at one loop.
